\documentclass[11pt]{article}
\usepackage[margin=1in]{geometry}
\usepackage{amsmath,amssymb,amsthm}
\usepackage{bm}
\usepackage[colorlinks=true,linkcolor=blue,citecolor=blue,urlcolor=blue]{hyperref}

\newcommand{\KL}{\mathrm{KL}}
\newcommand{\E}{\mathbb{E}}
\newcommand{\Var}{\mathrm{Var}}
\newcommand{\Cov}{\mathrm{Cov}}
\newcommand{\X}{\mathcal{X}}
\newcommand{\Sset}{\mathcal{S}}
\newcommand{\Aset}{\mathcal{A}}
\newcommand{\F}{\mathcal{F}}
\newcommand{\hs}{h^\star}
\DeclareMathOperator*{\argmax}{arg\,max}
\DeclareMathOperator*{\argmin}{arg\,min}
\DeclareMathOperator*{\osc}{osc}
\newcommand{\proj}{\Pi}
\newcommand{\lmin}{\lambda_{\min}}
\newcommand{\lmax}{\lambda_{\max}}

\theoremstyle{plain}
\newtheorem{theorem}{Theorem}
\newtheorem{proposition}{Proposition}
\newtheorem{lemma}{Lemma}
\newtheorem{corollary}{Corollary}
\theoremstyle{definition}
\newtheorem{assumption}{Assumption}
\theoremstyle{remark}
\newtheorem{remark}{Remark}

\title{\textbf{Idealised terminal guarantees for the adaptive dual $h$-transform}\\[2pt]
\large constraint satisfaction and fidelity under a fixed budget and one-step dual ascent}
\date{}

\begin{document}
\maketitle

\section{Setup and notation}
\label{sec:setup}

A masked discrete diffusion model is sampled in reverse along a grid
$1=t_0>t_1>\dots>t_N=0$. The state space $\X$ is finite, $x_0$ is the deterministic
fully-masked prior, and step $k\in\{0,\dots,N-1\}$ is the transition
$x_k\to x_{k+1}$ governed by the unconditional reverse kernel
$p^{(k)}(\cdot\mid x_k)\triangleq p_\theta(x_{k+1}=\cdot\mid x_k)$ (so the multiplier
$\lambda_k$ and the state $x_k$ it conditions on carry the same index, matching the
companion paper). Write $P_\theta$ for the unconditional path law and
$\F_k\triangleq\sigma(x_0,\dots,x_k)$ for the natural filtration; $\F_k$ is the
information available \emph{before} $x_{k+1}$ is drawn.

The constraint fixes $S\subseteq\X$; the terminal clean sample must satisfy
$Y_T\in S$. Let
\[
\hs(x,t)\triangleq P\big(Y_T\in S\mid Y_t=x\big)\in[0,1]
\]
be the true Doob $h$-function, with $\hs(\cdot,t_N)=\hs(\cdot,0)=\mathbf 1\{\cdot\in S\}$.
By the tower property it is a $P_\theta$-martingale: for every $k$ and $x_k$,
\begin{equation}
\label{eq:mart}
\sum_{y}p^{(k)}(y\mid x_k)\,\hs(y,t_{k+1})=\hs(x_k,t_k).
\end{equation}
To see this, condition on the next state $x_{k+1}$ and use that the reverse chain is
Markov, so that given $x_{k+1}$ the terminal event $\{Y_T\in S\}$ is independent of $x_k$:
\begin{align*}
\hs(x_k,t_k)
&=P\big(Y_T\in S\mid x_k\big)
=\E\big[\mathbf 1\{Y_T\in S\}\mid x_k\big]\\
&=\sum_{y}\underbrace{P\big(x_{k+1}=y\mid x_k\big)}_{=\,p^{(k)}(y\mid x_k)}\,
\E\big[\mathbf 1\{Y_T\in S\}\mid x_{k+1}=y,\,x_k\big]\\
&=\sum_{y}p^{(k)}(y\mid x_k)\,
\underbrace{\E\big[\mathbf 1\{Y_T\in S\}\mid x_{k+1}=y\big]}_{=\,\hs(y,t_{k+1})},
\end{align*}
where the second line is the tower property over $\F_{k+1}$ and the third uses the
Markov property to drop $x_k$. Equivalently,
$\E_\theta[\hs(x_{k+1},t_{k+1})\mid\F_k]=\hs(x_k,t_k)$, i.e.\ $\{\hs(x_k,t_k)\}_k$
is a $P_\theta$-martingale.

\paragraph{Tempered and exact transforms.} For $\lambda\ge0$ define the tilted
kernel (convention $0^\lambda=0$ for $\lambda>0$)
\begin{equation}
\label{eq:tilt}
q^{(k)}_\lambda(y\mid x_k)=\frac{p^{(k)}(y\mid x_k)\,\hs(y,t_{k+1})^{\lambda}}{Z_k(\lambda)},
\qquad Z_k(\lambda)=\sum_y p^{(k)}(y\mid x_k)\,\hs(y,t_{k+1})^{\lambda}.
\end{equation}
At $\lambda=1$, \eqref{eq:mart} gives $Z_k(1)=\hs(x_k,t_k)$, so
$q^{(k)}_1=p^{(k)}\hs(\cdot,t_{k+1})/\hs(x_k,t_k)$ is the exact Doob
$h$-transform; $P^\star$ denotes its path law and $p^{(k)}_\star\triangleq q^{(k)}_1$.

\paragraph{The realised algorithm (fixed budget, one-step ascent).} Fix a budget
$C\ge0$, a step size $\rho>0$, and a box $[\lmin,\lmax]$ with $0<\lmin\le1\le\lmax$.
Initialise $\lambda_0=1$. For $k=0,\dots,N-1$:
\begin{equation}
\label{eq:alg}
x_{k+1}\sim q^{(k)}_{\lambda_k}(\cdot\mid x_k),\qquad
\lambda_{k+1}=\proj_{[\lmin,\lmax]}\!\Big(\lambda_k-\rho\big(\log\hs(x_{k+1},t_{k+1})+C\big)\Big)\ \ (k<N-1).
\end{equation}
Each $\lambda_k$ is a deterministic function of $x_0,\dots,x_k$, hence
$\F_k$-measurable; the update $\lambda_k\!\to\!\lambda_{k+1}$ uses the fresh draw
$x_{k+1}$. Let $Q$ be the resulting path law, $q(x_N)$ its terminal marginal,
$A\triangleq\{Y_T\in S\}$, and $\pi_0\triangleq P_\theta(A)\in(0,1]$.

\paragraph{Goal.} Bound $\KL\big(q(x_N)\,\|\,p_{\text{data}}(\cdot\mid S)\big)$ and
$P_Q(x_N\in S)$, where $p_{\text{data}}(\cdot\mid S)$ is the data law restricted to
$S$ (Prop.~\ref{prop:iproj}).

\section{Assumptions}
\label{sec:ass}

\begin{assumption}[idealised regime]
\label{as:exact}
\textup{(a)} the discriminator is exact, $h_\phi=\hs$ in \eqref{eq:tilt}--\eqref{eq:alg};
\textup{(b)} the base model is exact, $p_\theta(x_N=\cdot)=p_{\text{data}}(\cdot)$.
\end{assumption}

\begin{assumption}[alive support and bounded range]
\label{as:alive}
For a state $x$ at step $k$ let $\Aset_k(x)\triangleq\{y:p^{(k)}(y\mid x)>0,\ \hs(y,t_{k+1})>0\}$
be the alive successor support. There is $\epsilon\in(0,1]$, uniform in $k,N$, with
$\hs\ge\epsilon$ on every $\Aset_k(x_k)$ swept by $Q$ for $k\le N-2$, so
$R\triangleq\log(1/\epsilon)$ bounds $\osc_{\Aset_k}(\log\hs(\cdot,t_{k+1}))$. At the terminal step
$k=N-1$, $\hs(\cdot,t_N)=\mathbf 1\{\cdot\in S\}$, so $\hs\equiv1$ on $\Aset_{N-1}$ (no oscillation).
\end{assumption}

\begin{assumption}[nondegeneracy on the pre-terminal steps]
\label{as:mu}
There are $0<\mu\le L\le R^2/4$, uniform in $k,N$, with
$\Var_{q^{(k)}_\lambda}(\log\hs(\cdot,t_{k+1}))\in[\mu,L]$ for all $\lambda\in[\lmin,\lmax]$,
all $x_k$, and all $k\le N-2$. (At $k=N-1$ the variance is $0$; that step is
excluded from every $\mu$-dependent statement below.)
\end{assumption}

\begin{assumption}[time-regularity of the per-step program]
\label{as:drift}
With $dt_k\triangleq t_k-t_{k+1}>0$, there is $G<\infty$, uniform in $k,N$, such that
$\sup_{\lambda\in[\lmin,\lmax]}\big|\,\E_{q^{(k+1)}_\lambda}[-\log\hs(\cdot,t_{k+2})]
-\E_{q^{(k)}_\lambda}[-\log\hs(\cdot,t_{k+1})]\,\big|\le G\,dt_k$ almost surely along $Q$,
for $k\le N-3$.
\end{assumption}

\section{The gold target is the $I$-projection}
\label{sec:target}

\begin{proposition}
\label{prop:iproj}
Among path measures $Q$ with $Q(A)=1$ and $Q\ll P_\theta$ (else $\KL(Q\|P_\theta)=\infty$),
the exact transform $P^\star=P_\theta(\cdot\mid A)$ uniquely minimises $\KL(Q\|P_\theta)$,
with $\KL(P^\star\|P_\theta)=-\log\pi_0$, and its terminal marginal is the
constrained data law $p_{\text{data}}(\cdot\mid S)$.
\end{proposition}
\begin{proof}
Let $Q(A)=1$ and $Q\ll P_\theta$. Using $P_\theta=\pi_0P^\star$ on $A$, and that $Q$ charges only $A$,
\begin{gather*}
\KL(Q\|P_\theta)=\E_Q\!\Big[\log\tfrac{dQ}{dP_\theta}\Big]
=\E_Q\!\Big[\log\tfrac{dQ}{dP^\star}\Big]-\log\pi_0
=\KL(Q\|P^\star)-\log\pi_0\ \ge\ -\log\pi_0,
\end{gather*}
with equality iff $\KL(Q\|P^\star)=0$, i.e.\ $Q=P^\star$. For the terminal marginal,
\begin{gather*}
P^\star(x_N=z)=\frac{p_\theta(x_N=z)\,\mathbf 1\{z\in S\}}{\pi_0}
=\frac{p_{\text{data}}(z)\,\mathbf 1\{z\in S\}}{p_{\text{data}}(S)}=p_{\text{data}}(z\mid S),
\end{gather*}
using Assumption~\ref{as:exact}(b) and $\pi_0=p_{\text{data}}(S)$.
\end{proof}

Thus ``satisfy the constraint while staying closest to the data'' is exactly
``approximate $P^\star$''; all terminal bounds below measure the gap of $q(x_N)$
to $p_{\text{data}}(\cdot\mid S)$.

\section{Constraint satisfaction is structural}
\label{sec:sat}

\begin{theorem}
\label{thm:sat}
Under Assumption~\ref{as:exact}:
\textup{(a)} if $\lambda_k\equiv1$ then $Q=P^\star$, so $P_Q(x_N\in S)=1$ and
$q(x_N)=p_{\text{data}}(\cdot\mid S)$ ($\KL=0$);
\textup{(b)} for \emph{any} schedule with $\lambda_k\ge\lmin>0$ for all $k$ --- in
particular the algorithm \eqref{eq:alg} --- one has $P_Q(x_N\in S)=1$.
\end{theorem}
\begin{proof}
\textup{(a)} With $\lambda_k\equiv1$ each kernel is $q^{(k)}_1=p^{(k)}\hs(\cdot,t_{k+1})/\hs(x_k,t_k)$
(a valid kernel since $Z_k(1)=\hs(x_k,t_k)$ by~\eqref{eq:mart}). The path
probability telescopes ($x_0$ deterministic):
\begin{gather*}
Q(x_1,\dots,x_N)=\prod_{k=0}^{N-1} q^{(k)}_1(x_{k+1}\mid x_k)
=\Big(\textstyle\prod_{k=0}^{N-1} p^{(k)}(x_{k+1}\mid x_k)\Big)\prod_{k=0}^{N-1}\frac{\hs(x_{k+1},t_{k+1})}{\hs(x_k,t_k)},\\
\prod_{k=0}^{N-1}\frac{\hs(x_{k+1},t_{k+1})}{\hs(x_k,t_k)}
=\frac{\hs(x_N,t_N)}{\hs(x_0,t_0)}=\frac{\mathbf 1\{x_N\in S\}}{\pi_0}
\quad\big(\hs(\cdot,t_N)=\mathbf 1\{\cdot\in S\},\ \hs(x_0,t_0)=\pi_0\big),\\
\Longrightarrow\quad
Q(x_1,\dots,x_N)=\underbrace{\Big(\textstyle\prod_{k=0}^{N-1} p^{(k)}(x_{k+1}\mid x_k)\Big)}_{=\,P_\theta(x_1,\dots,x_N)}
\,\frac{\mathbf 1\{x_N\in S\}}{\pi_0}=P_\theta(\cdot\mid A)=P^\star,
\end{gather*}
the middle $\hs(x_1,t_1),\dots,\hs(x_{N-1},t_{N-1})$ cancelling in pairs. Conclude by Prop.~\ref{prop:iproj}.

\textup{(b)} By induction on $k$ we show $\hs(x_k,t_k)>0$ $Q$-a.s.:
\begin{gather*}
\hs(x_0,t_0)=\pi_0>0;\\
\hs(x_k,t_k)>0\ \Longrightarrow\ \textstyle\sum_y p^{(k)}(y\mid x_k)\,\hs(y,t_{k+1})=\hs(x_k,t_k)>0,\\
q^{(k)}_{\lambda_k}(\cdot\mid x_k)\propto p^{(k)}(\cdot\mid x_k)\,\hs(\cdot,t_{k+1})^{\lambda_k}\ \ (\lambda_k\ge\lmin>0)
\ \Longrightarrow\ \hs(x_{k+1},t_{k+1})>0\ \text{a.s.}
\end{gather*}
At the terminal step, $\hs(x_N,t_N)=\mathbf 1\{x_N\in S\}$, so $\hs(x_N,t_N)>0\Rightarrow x_N\in S$, i.e.\ $P_Q(x_N\in S)=1$.
\end{proof}

In the exact regime the hard constraint holds with probability one for any positive
schedule; the budget $C$ and the one-step ascent affect only fidelity.

\section{Master fidelity bound}
\label{sec:master}

\begin{lemma}[KL between two tilts of a common base]
\label{lem:osc}
Let $q_i\propto p\,e^{U_i}$ ($i=1,2$) share base $p$, set $\Delta=U_1-U_2$ and
$\omega=\osc_{\operatorname{supp}p}(\Delta)$. If $\omega<\infty$ then $\KL(q_1\|q_2)\le\omega^2/8$.
\end{lemma}
\begin{proof}
Write $q_i(x)=p(x)e^{U_i(x)}/Z_i$. Pointwise, then summing,
\begin{gather*}
q_2(x)\,e^{\Delta(x)}=\frac{p(x)e^{U_1(x)}}{Z_2}=\frac{Z_1}{Z_2}\,q_1(x)
\ \Longrightarrow\
q_1(x)=\frac{q_2(x)e^{\Delta(x)}}{M(1)},\quad M(\tau)\triangleq\E_{q_2}[e^{\tau\Delta}],\ \ M(1)=\tfrac{Z_1}{Z_2}.
\end{gather*}
Set $q_\tau(x)=q_2(x)e^{\tau\Delta(x)}/M(\tau)$ (so $q_0=q_2$, and $q_1$ as above) and $\psi(\tau)=\log M(\tau)$. Differentiating the finite sum,
\begin{gather*}
\psi(0)=0,\qquad
\psi'(\tau)=\E_{q_\tau}[\Delta],\qquad
\psi''(\tau)=\Var_{q_\tau}(\Delta)\ge0.
\end{gather*}
Since $\log\tfrac{q_1}{q_2}=\Delta-\psi(1)$,
\begin{gather*}
\KL(q_1\|q_2)=\E_{q_1}[\Delta]-\psi(1)=\psi'(1)-\psi(1).
\end{gather*}
Using $\psi(0)=0$, the fundamental theorem of calculus, and Fubini on $\{0\le s\le u\le1\}$,
\begin{gather*}
\psi'(1)-\psi(1)=\int_0^1\!\big[\psi'(1)-\psi'(s)\big]\,ds
=\int_0^1\!\!\int_s^1\psi''(u)\,du\,ds
=\int_0^1 u\,\psi''(u)\,du.
\end{gather*}
On $\operatorname{supp}p=\operatorname{supp}q_u$, $\Delta$ ranges in an interval of length $\omega$, so Popoviciu gives $\psi''(u)=\Var_{q_u}(\Delta)\le\omega^2/4$, whence
\begin{gather*}
\KL(q_1\|q_2)=\int_0^1 u\,\psi''(u)\,du\ \le\ \frac{\omega^2}{4}\int_0^1 u\,du=\frac{\omega^2}{8}.
\end{gather*}
\end{proof}

\begin{theorem}[terminal fidelity is governed by $\sum_{k\le N-2}(\lambda_k-1)^2$]
\label{thm:master}
Under Assumptions~\ref{as:exact}--\ref{as:alive}, for any predictable schedule
$(\lambda_k)$ with $\lambda_k\in[\lmin,\lmax]$ ($\lmin>0$), writing
$\varepsilon_k=(\lambda_k-1)\log\hs(\cdot,t_{k+1})$,
\begin{equation}
\label{eq:master}
\KL\!\big(q(x_N)\,\big\|\,p_{\text{data}}(\cdot\mid S)\big)
\le\KL(Q\|P^\star)
=\sum_{k=0}^{N-1}\E_Q\big[\KL\big(q^{(k)}_{\lambda_k}\,\|\,p^{(k)}_\star\big)\big]
\le\frac18\sum_{k=0}^{N-2}\E_Q\big[\osc_{\Aset_k}(\varepsilon_k)^2\big]
\le\frac{R^2}{8}\sum_{k=0}^{N-2}\E_Q\big[(\lambda_k-1)^2\big].
\end{equation}
\end{theorem}
\begin{proof}
For $\lambda_k>0$ both $q^{(k)}_{\lambda_k}$ and $p^{(k)}_\star=q^{(k)}_1$ are supported on
$\{p^{(k)}>0,\hs(\cdot,t_{k+1})>0\}$, hence mutually absolutely continuous with finite per-step $\KL$, and $Q\ll P^\star\ll P_\theta$. Then
\begin{gather*}
\KL\!\big(q(x_N)\,\|\,p_{\text{data}}(\cdot\mid S)\big)
\ \overset{\text{DPI}}{\le}\ \KL(Q\|P^\star)
\ \overset{\text{chain rule}}{=}\ \sum_{k=0}^{N-1}\E_Q\!\big[\KL\big(Q(x_{k+1}\mid\F_k)\,\|\,p^{(k)}_\star(\cdot\mid x_k)\big)\big],\\
Q(x_{k+1}\mid\F_k)=q^{(k)}_{\lambda_k}\quad(\lambda_k\ \text{is}\ \F_k\text{-measurable}).
\end{gather*}
Each kernel is a tilt of $p^{(k)}(\cdot\mid x_k)$ with $\Delta=\varepsilon_k$ on $\Aset_k$, so by Lemma~\ref{lem:osc} and Assumption~\ref{as:alive},
\begin{gather*}
\KL\big(q^{(k)}_{\lambda_k}\,\|\,p^{(k)}_\star\big)\le\tfrac18\osc_{\Aset_k}(\varepsilon_k)^2,\qquad
\osc_{\Aset_k}(\varepsilon_k)=|\lambda_k-1|\,\osc_{\Aset_k}(\log\hs(\cdot,t_{k+1}))\le|\lambda_k-1|\,R.
\end{gather*}
At the terminal step $k=N-1$, $\hs\equiv1$ on $\Aset_{N-1}$, so $\varepsilon_{N-1}\equiv0$ and that term vanishes; summing over $k\le N-2$ gives \eqref{eq:master}.
\end{proof}

The entire fidelity question is therefore: how far does the realised $\lambda_k$
($k\le N-2$) stray from the Doob exponent $1$?

\section{The per-step dual and the calibration gap}
\label{sec:dual}

For $k\le N-2$ the update \eqref{eq:alg} is projected ascent on the concave dual of
$\min_q\KL(q\|p^{(k)})$ s.t.\ $\E_q[-\log\hs(\cdot,t_{k+1})]\le C$ (given $x_k$). With
$S_k(\lambda)\triangleq\E_{q^{(k)}_\lambda}[-\log\hs(\cdot,t_{k+1})]$, the dual is
$g^{(k)}(\lambda)=-\log Z_k(\lambda)-\lambda C$ with
\begin{equation}
\label{eq:gder}
g^{(k)}{}'(\lambda)=S_k(\lambda)-C,\qquad
g^{(k)}{}''(\lambda)=-\Var_{q^{(k)}_\lambda}(\log\hs(\cdot,t_{k+1}))\in[-L,-\mu],
\end{equation}
so $S_k$ is strictly decreasing with $S_k'\in[-L,-\mu]$ (Assumption~\ref{as:mu}).
Since $g^{(k)}$ is continuous and strictly concave on the compact box $[\lmin,\lmax]$,
the maximiser $\lambda^\star_k(x_k;C)\triangleq\argmax_{[\lmin,\lmax]}g^{(k)}$ exists
and is unique; it is the attractor of \eqref{eq:alg}. Let
\begin{equation}
\label{eq:Cstar}
C^\star_k(x_k)\triangleq S_k(1)=\E_{q^{(k)}_1}[-\log\hs(\cdot,t_{k+1})\mid x_k]
\end{equation}
be the (state-dependent) surprisal the exact transform spends at step $k$.

\begin{lemma}[calibration gap]
\label{lem:calib}
For $k\le N-2$, with $\lmin\le1\le\lmax$, the dual maximiser obeys
\[
b_k\triangleq|\lambda^\star_k-1|\;\le\;\frac{|C-C^\star_k(x_k)|}{\mu}.
\]
\end{lemma}
\begin{proof}
$S_k$ is strictly decreasing with $|S_k'|\ge\mu$ and $S_k(1)=C^\star_k$. Three cases for $\lambda^\star_k$:
\begin{gather*}
\text{interior:}\quad S_k(\lambda^\star_k)=C\ \Rightarrow\ C-C^\star_k=S_k'(\xi)(\lambda^\star_k-1)\ \Rightarrow\ b_k\le\tfrac{|C-C^\star_k|}{\mu};\\
\lambda^\star_k=\lmax\ge1:\quad g^{(k)}{}'(\lmax)\ge0\Rightarrow S_k(\lmax)\ge C,\ \ S_k(\lmax)\le S_k(1)=C^\star_k
\ \Rightarrow\ C^\star_k-C\ge\mu\,b_k;\\
\lambda^\star_k=\lmin\le1:\quad g^{(k)}{}'(\lmin)\le0\Rightarrow S_k(\lmin)\le C,\ \ S_k(\lmin)\ge S_k(1)=C^\star_k
\ \Rightarrow\ C-C^\star_k\ge\mu\,b_k.
\end{gather*}
Each case uses $|S_k'(\xi)|\ge\mu$ and $b_k=|\lambda^\star_k-1|$; in all of them $b_k\le|C-C^\star_k|/\mu$.
\end{proof}

Since $C^\star_k$ varies with $k$ (typically decreasing as $t\to0$, where $\hs\to1$ on
the conditioned path so $-\log\hs\to0$), a single fixed $C$ cannot make
$\lambda^\star_k=1$ at all steps; that requires the calibrated step-dependent budget
$C=C^\star_k$.

\section{One-step stochastic tracking}
\label{sec:track}

Let $a_k\triangleq\lambda_k-\lambda^\star_k$, $\delta_k\triangleq|\lambda^\star_{k+1}-\lambda^\star_k|$
(defined for $k\le N-3$, so both optima are pre-terminal),
$\widehat{g'}_k\triangleq-(\log\hs(x_{k+1},t_{k+1})+C)$, and $\zeta_k\triangleq\widehat{g'}_k-g^{(k)}{}'(\lambda_k)$.

\begin{remark}[the multiplier lags one step]
\label{rem:lag}
The algorithm \eqref{eq:alg} samples \emph{before} it updates: step $k$ is tilted by $\lambda_k$
(conditioned on $x_k$), and the fresh draw $x_{k+1}$ then forms $\lambda_{k+1}$ for the \emph{next}
step. Thus $\lambda_k$ is a \emph{one-step-lagged} tracker --- it is a projected-ascent step on
the step-$(k{-}1)$ dual $g^{(k-1)}$ but is applied to step $k$, whose own optimum $\lambda^\star_k$
it should match. This lag is exactly the drift term in the recursion below: $\lambda_{k+1}$
contracts toward $\lambda^\star_k$ (the dual it ascends) and is then compared against
$\lambda^\star_{k+1}$ (the step it tilts), leaving the per-step gap
$\delta_k=|\lambda^\star_{k+1}-\lambda^\star_k|$. The unbiasedness
$\E[\widehat{g'}_k\mid\F_k]=g^{(k)}{}'(\lambda_k)$ (Lemma~\ref{lem:noise}(a)) holds precisely
because $x_{k+1}\sim q^{(k)}_{\lambda_k}$ is drawn at the same $\F_k$-measurable $\lambda_k$ that
tilts the step. The lag-free variant --- solving each frozen-step dual to $\lambda^\star_k$ before
drawing $x_{k+1}$ --- is the $\delta_k\to0$ limit (equivalently $\overline{dt}\to0$ in Thm.~\ref{thm:track}).
\end{remark}

\begin{lemma}[unbiased, bounded-variance gradient; drift stability]
\label{lem:noise}
For $k\le N-2$: \textup{(a)} since $x_{k+1}\sim q^{(k)}_{\lambda_k}(\cdot\mid x_k)$ and $\lambda_k$
is $\F_k$-measurable, $\E[\zeta_k\mid\F_k]=0$ and $\E[\zeta_k^2\mid\F_k]
=\Var_{q^{(k)}_{\lambda_k}}(\log\hs(\cdot,t_{k+1}))\le R^2/4=:\sigma_g^2$. \textup{(b)} for $k\le N-3$,
$\delta_k\le G\,dt_k/\mu$ a.s.
\end{lemma}
\begin{proof}
\textup{(a)} Since $x_{k+1}\sim q^{(k)}_{\lambda_k}$ given $\F_k$,
\begin{gather*}
\E[\widehat{g'}_k\mid\F_k]=\E_{q^{(k)}_{\lambda_k}}[-\log\hs(\cdot,t_{k+1})]-C=g^{(k)}{}'(\lambda_k)
\ \Rightarrow\ \E[\zeta_k\mid\F_k]=0,\\
\E[\zeta_k^2\mid\F_k]=\Var_{q^{(k)}_{\lambda_k}}(\log\hs(\cdot,t_{k+1}))\le R^2/4=\sigma_g^2\quad(\text{Popoviciu},\ \log\hs\in[\log\epsilon,0]).
\end{gather*}
\textup{(b)} Write $\Delta\triangleq\lambda^\star_{k+1}-\lambda^\star_k$ (so $\delta_k=|\Delta|$). First-order optimality of the concave $g^{(k)},g^{(k+1)}$ on $[\lmin,\lmax]$, each tested at the other step's maximiser, and their difference give
\begin{gather*}
g^{(k)}{}'(\lambda^\star_k)\,\Delta\le0\le g^{(k+1)}{}'(\lambda^\star_{k+1})\,\Delta
\ \Longrightarrow\ \big[g^{(k+1)}{}'(\lambda^\star_{k+1})-g^{(k)}{}'(\lambda^\star_k)\big]\Delta\ge0.
\end{gather*}
Split through $g^{(k+1)}{}'(\lambda^\star_k)$ and use $\mu$-strong concavity of $g^{(k+1)}$ (so $(g^{(k+1)}{}'(a)-g^{(k+1)}{}'(b))(a-b)\le-\mu(a-b)^2$):
\begin{gather*}
\underbrace{\big[g^{(k+1)}{}'(\lambda^\star_{k+1})-g^{(k+1)}{}'(\lambda^\star_k)\big]\Delta}_{\le\,-\mu\Delta^2}
+\big[g^{(k+1)}{}'(\lambda^\star_k)-g^{(k)}{}'(\lambda^\star_k)\big]\Delta\ \ge\ 0,\\
\Longrightarrow\quad
\mu\Delta^2\ \le\ \big[g^{(k+1)}{}'(\lambda^\star_k)-g^{(k)}{}'(\lambda^\star_k)\big]\Delta
\ \le\ \sup_{\lambda}\big|g^{(k+1)}{}'(\lambda)-g^{(k)}{}'(\lambda)\big|\,|\Delta|.
\end{gather*}
If $\Delta=0$ the bound is trivial; otherwise divide by $|\Delta|=\delta_k$. Finally
$g^{(k)}{}'(\lambda)=\E_{q^{(k)}_\lambda}[-\log\hs(\cdot,t_{k+1})]-C$ by~\eqref{eq:gder} (the $-C$ cancels in the difference), so the supremum is exactly the quantity bounded in Assumption~\ref{as:drift}:
\begin{gather*}
\mu\,\delta_k\ \le\ \sup_{\lambda}\big|g^{(k+1)}{}'(\lambda)-g^{(k)}{}'(\lambda)\big|\ \le\ G\,dt_k.
\end{gather*}
\end{proof}

\begin{lemma}[per-step contraction]
\label{lem:contr}
For $0<\rho\le1/L$ and $k\le N-2$, $T_k(\lambda)=\proj_{[\lmin,\lmax]}(\lambda+\rho g^{(k)}{}'(\lambda))$
fixes $\lambda^\star_k$ and is $(1-\rho\mu)$-Lipschitz.
\end{lemma}
\begin{proof}
$\lambda^\star_k$ solves the projected fixed-point equation $\lambda=T_k(\lambda)$, and
\begin{gather*}
\frac{d}{d\lambda}\big(\lambda+\rho g^{(k)}{}'(\lambda)\big)=1-\rho\Var_{q^{(k)}_\lambda}(\log\hs(\cdot,t_{k+1}))\in[1-\rho L,\,1-\rho\mu]\subseteq[0,1-\rho\mu]\quad(\rho\le1/L),
\end{gather*}
so the inner map is $(1-\rho\mu)$-Lipschitz; composing with the nonexpansive $\proj$ preserves this.
\end{proof}

\begin{theorem}[summed second-moment tracking]
\label{thm:track}
Under Assumptions~\ref{as:exact}--\ref{as:drift} and $0<\rho\le1/L$ (so $\rho\mu\le1$),
\begin{equation}
\label{eq:track}
\sum_{k=0}^{N-2}\E[a_k^2]\;\le\;\frac{2\,\E[a_0^2]}{\rho\mu}
\;+\;\frac{3(N-1)\,\rho\,\sigma_g^2}{\mu}
\;+\;\frac{6\,G^2\,\overline{dt}}{\rho^2\mu^4},
\qquad \overline{dt}\triangleq\max_k dt_k,
\end{equation}
where $a_0=1-\lambda^\star_0$, so $\E[a_0^2]=\E[b_0^2]\le\mu^{-2}\E[(C-C^\star_0)^2]$.
\end{theorem}
\begin{proof}
Fix $k\le N-3$; we derive a one-step recursion for $\E[a_{k+1}^2]$ in five steps and unroll.
Throughout $\beta\triangleq\rho\mu/2$ and $\rho\mu\le1$ (since $\rho\le1/L$ and $\mu\le L$).

\emph{Step 1 (the deterministic part contracts).} Let $S(\lambda)\triangleq\lambda+\rho g^{(k)}{}'(\lambda)$
be the inner ascent map ($T_k=\proj\circ S$), and $w_k\triangleq S(\lambda_k)-S(\lambda^\star_k)$.
Lemma~\ref{lem:contr} ($S$ is $(1-\rho\mu)$-Lipschitz) gives
\begin{gather*}
|w_k|\le(1-\rho\mu)|a_k|,\qquad w_k\ \text{is}\ \F_k\text{-measurable}.
\end{gather*}

\emph{Step 2 (one noisy step, conditional second moment).} The update is
$\lambda_{k+1}=\proj(S(\lambda_k)+\rho\zeta_k)$ and the optimum is the fixed point
$\lambda^\star_k=\proj(S(\lambda^\star_k))$; nonexpansiveness of $\proj$ gives
$|\lambda_{k+1}-\lambda^\star_k|\le|w_k+\rho\zeta_k|$. Squaring, taking $\E[\cdot\mid\F_k]$,
and using $\E[\zeta_k\mid\F_k]=0$, $\E[\zeta_k^2\mid\F_k]\le\sigma_g^2$ (Lemma~\ref{lem:noise}(a)),
\begin{gather*}
\E\big[(\lambda_{k+1}-\lambda^\star_k)^2\mid\F_k\big]
\le w_k^2+2\rho\,w_k\underbrace{\E[\zeta_k\mid\F_k]}_{=\,0}+\rho^2\underbrace{\E[\zeta_k^2\mid\F_k]}_{\le\,\sigma_g^2}
\le(1-\rho\mu)^2a_k^2+\rho^2\sigma_g^2.
\end{gather*}

\emph{Step 3 (absorb the moving target, Young).} Since
$a_{k+1}=(\lambda_{k+1}-\lambda^\star_k)+(\lambda^\star_k-\lambda^\star_{k+1})$ with
$|\lambda^\star_k-\lambda^\star_{k+1}|=\delta_k$, Young's inequality
$(x+y)^2\le(1+\beta)x^2+(1+\beta^{-1})y^2$, then Step 2 and $\delta_k\le G\,dt_k/\mu$ a.s.\ (Lemma~\ref{lem:noise}(b)), give
\begin{gather*}
a_{k+1}^2\le(1+\beta)(\lambda_{k+1}-\lambda^\star_k)^2+(1+\beta^{-1})\delta_k^2,\\
\E[a_{k+1}^2\mid\F_k]\le(1+\beta)(1-\rho\mu)^2a_k^2+(1+\beta)\rho^2\sigma_g^2+(1+\beta^{-1})\frac{G^2dt_k^2}{\mu^2}.
\end{gather*}

\emph{Step 4 (the three constants, via $x\triangleq\rho\mu\in(0,1]$).}
\begin{gather*}
(1+\beta)(1-\rho\mu)^2-(1-\beta)=\big(1+\tfrac x2\big)(1-x)^2-\big(1-\tfrac x2\big)=x\big(\tfrac{x^2}{2}-1\big)\le0,\\
1+\beta=1+\tfrac x2\le\tfrac32,\qquad 1+\beta^{-1}=1+\tfrac2x\le\tfrac3x=\tfrac{3}{\rho\mu},
\end{gather*}
so $(1+\beta)(1-\rho\mu)^2\le1-\beta=1-\tfrac{\rho\mu}{2}$. Substituting into Step 3 and taking full expectation,
\begin{gather*}
\E[a_{k+1}^2]\le\big(1-\tfrac{\rho\mu}{2}\big)\E[a_k^2]+\tfrac32\rho^2\sigma_g^2+\tfrac{3}{\rho\mu}\frac{G^2dt_k^2}{\mu^2}.
\end{gather*}

\emph{Step 5 (unroll and sum).} The recursion is affine with contraction factor $1-\tfrac{\rho\mu}{2}$;
unrolling it for $k=0,\dots,N-3$ (with $\sum_{j\ge0}(1-\tfrac{\rho\mu}{2})^j=\tfrac{2}{\rho\mu}$), summing the
resulting $\E[a_k^2]$ over $k=0,\dots,N-2$, and using
$\sum_k dt_k^2\le\overline{dt}\sum_k dt_k\le\overline{dt}$,
\begin{gather*}
\sum_{k=0}^{N-2}\E[a_k^2]\le\frac{2\,\E[a_0^2]}{\rho\mu}
+\frac{2}{\rho\mu}\Big(\tfrac32(N-1)\rho^2\sigma_g^2+\tfrac{3}{\rho\mu}\frac{G^2}{\mu^2}\overline{dt}\Big)
=\frac{2\,\E[a_0^2]}{\rho\mu}+\frac{3(N-1)\rho\sigma_g^2}{\mu}+\frac{6G^2\overline{dt}}{\rho^2\mu^4},\\
a_0=1-\lambda^\star_0\ \Rightarrow\ \E[a_0^2]=\E[b_0^2]\le\mu^{-2}\E[(C-C^\star_0)^2]\quad(\text{Lemma~\ref{lem:calib}}).
\end{gather*}
\end{proof}

\section{Main result}
\label{sec:main}

\begin{theorem}[exact regime, fixed $C$, one-step ascent]
\label{thm:main}
Under Assumptions~\ref{as:exact}--\ref{as:drift}, run \eqref{eq:alg} with $0<\rho\le1/L$
and $\lambda_0=1$. Then $P_Q(x_N\in S)=1$, and
\begin{equation}
\label{eq:main}
\KL\!\big(q(x_N)\,\big\|\,p_{\text{data}}(\cdot\mid S)\big)
\;\le\;\frac{R^2}{4}\underbrace{\sum_{k=0}^{N-2}\E_Q[a_k^2]}_{\textnormal{(B) one-step tracking}}
\;+\;\frac{R^2}{4}\underbrace{\sum_{k=0}^{N-2}\E_Q[b_k^2]}_{\textnormal{(A) fixed-}C\textnormal{ calibration}},
\end{equation}
where
\[
\textnormal{(A)}\le\frac{1}{\mu^2}\sum_{k=0}^{N-2}\E_Q\big[(C-C^\star_k(x_k))^2\big],
\qquad
\textnormal{(B)}\le\frac{2\,\E[b_0^2]}{\rho\mu}+\frac{3(N-1)\rho\sigma_g^2}{\mu}+\frac{6G^2\overline{dt}}{\rho^2\mu^4}.
\]
\end{theorem}
\begin{proof}
Satisfaction is Thm.~\ref{thm:sat}(b) ($\lambda_k\ge\lmin>0$). For fidelity, $(\lambda_k-1)^2\le2a_k^2+2b_k^2$ in Thm.~\ref{thm:master} gives
\begin{gather*}
\KL\!\big(q(x_N)\,\|\,p_{\text{data}}(\cdot\mid S)\big)\le\frac{R^2}{8}\sum_{k=0}^{N-2}\E_Q\big[2a_k^2+2b_k^2\big]
=\frac{R^2}{4}\Big(\sum_{k=0}^{N-2}\E_Q[a_k^2]+\sum_{k=0}^{N-2}\E_Q[b_k^2]\Big);
\end{gather*}
substitute Lemma~\ref{lem:calib} ($b_k\le|C-C^\star_k|/\mu$, then $\E_Q$) for \textnormal{(A)} and Theorem~\ref{thm:track} for \textnormal{(B)}.
\end{proof}

\begin{corollary}[best constant budget, deterministic profile]
\label{cor:bestC}
Suppose the surprisal profile is state-independent, $C^\star_k(x_k)\equiv c^\star_k$
(depends on $t_{k+1}$ only). Then over fixed budgets $C$, term \textnormal{(A)} is minimised
at $C=\bar c^\star=\frac{1}{N-1}\sum_{k=0}^{N-2}c^\star_k$, with
\[
\min_C\ \textnormal{(A)}=\frac{1}{\mu^2}\sum_{k=0}^{N-2}(c^\star_k-\bar c^\star)^2
=\frac{N-1}{\mu^2}\,\mathrm{Var}_k(c^\star_k).
\]
\end{corollary}
\begin{proof}
With $c^\star_k$ deterministic, $\textnormal{(A)}=\mu^{-2}\sum_{k=0}^{N-2}(C-c^\star_k)^2$ is quadratic in $C$, so
\begin{gather*}
\argmin_C\sum_{k=0}^{N-2}(C-c^\star_k)^2=\bar c^\star,\qquad
\min_C\sum_{k=0}^{N-2}(C-c^\star_k)^2=\sum_{k=0}^{N-2}(c^\star_k-\bar c^\star)^2=(N-1)\,\mathrm{Var}_k(c^\star_k).
\end{gather*}
\end{proof}

\begin{remark}[reading the bound]
\label{rem:read}
\textup{(i)} \emph{Satisfaction is free}: $P_Q(x_N\in S)=1$ for any fixed $C$ and any
$\rho$; the budget and one-step ascent cost only fidelity.
\textup{(ii)} \emph{(A) is the price of a fixed budget}: $\rho$-independent, vanishing
only under the calibrated step-dependent budget $C=C^\star_k$ (i.e.\ $\lambda\equiv1$,
the exact transform); for a fixed $C$ it scales with the cross-step spread of $C^\star_k$
(Cor.~\ref{cor:bestC}).
\textup{(iii)} \emph{(B) is the price of one-step ascent}: a single projected step per
reverse step leaves a single-sample noise term $\propto N\rho\sigma_g^2/\mu$ and a
moving-target drift term $\propto G^2\overline{dt}/(\rho^2\mu^4)$ (no $N$ factor; vanishes
as the grid refines), traded by $\rho$.
\textup{(iv)} \emph{Scaling}: with $R,\mu,L,G$ uniform in $N$, the drift part of (B)
$\to0$ as $\overline{dt}\to0$, while the noise part $N\rho\sigma_g^2/\mu$ and (A) grow
linearly in $N$ unless $\rho$ and the $C$-mismatch shrink; the bound is informative when
these per-step floors are small and becomes vacuous ($\Theta(N)$) for a badly mismatched
$C$ or too large $\rho$. The constants blow up if $\mu\to0$ (near-deterministic late steps),
where the bounded-oscillation control of those steps should be used directly instead.
\end{remark}


% ===========================================================================
\section{Imperfect discriminator ($\eta>0$): why adaptive tempering is needed}
\label{sec:eta}

Sections~\ref{sec:target}--\ref{sec:main} assumed $h_\phi=\hs$, under which $\lambda\equiv1$
is optimal and adaptivity is unnecessary. We now drop that: the \emph{realised} algorithm
runs \eqref{eq:alg} with the \emph{learned} $h_\phi$ in place of $\hs$, i.e.\ its step kernel
is $q^{(k)}_\lambda[h_\phi](y\mid x_k)\propto p^{(k)}(y\mid x_k)\,h_\phi(y,t_{k+1})^{\lambda}$,
while the gold target is unchanged, $P^\star=P_\theta(\cdot\mid A)$ with terminal
$p_{\text{data}}(\cdot\mid S)$ (built from the \emph{true} $\hs$); the base model stays exact
(Assumption~\ref{as:exact}(b)). Two error types now coexist, and we keep them apart:
\emph{calibration} error on states both discriminators call alive, and \emph{support mismatch}.

\begin{assumption}[imperfect discriminator: common support and false positives]
\label{as:eta}
For a state $x_k$ at step $k$ split the $h_\phi$-positive base support
$\{y:p^{(k)}(y\mid x_k)>0,\ h_\phi(y,t_{k+1})>0\}$ into the \emph{common alive support}
$\Aset^\cap_k(x_k)\triangleq\{y:h_\phi(y,t_{k+1})>0,\ \hs(y,t_{k+1})>0\}$ and the
\emph{false-positive set} $\Aset^+_k(x_k)\triangleq\{y:h_\phi(y,t_{k+1})>0,\ \hs(y,t_{k+1})=0\}$.
On every $\Aset^\cap_k$ both $\log h_\phi,\log\hs\in[\log\epsilon,0]$, and
\[
\eta\triangleq\sup_{0\le k\le N-1}\ \sup_{x_k}\ \sup_{y\in\Aset^\cap_k(x_k)}
\big|\log h_\phi(y,t_{k+1})-\log\hs(y,t_{k+1})\big|<\infty
\]
(no finiteness is claimed on $\Aset^+_k$, where $\log\hs=-\infty$). At the terminal step $k=N-1$
the common support is $\{y\in S:p^{(N-1)}(y\mid x_{N-1})>0\}$, where $\hs=1$, so there the error is
$-\log h_\phi\in[0,\log(1/\epsilon)]$ and is covered by $\eta$.
\end{assumption}

\begin{assumption}[regularity of the learned dual]
\label{as:eta-reg}
For $k\le N-2$ and $\lambda\in[\lmin,\lmax]$, on the realised support
$\{y:p^{(k)}(y\mid x_k)>0,\ h_\phi(y,t_{k+1})>0\}$: $\osc(\log h_\phi(\cdot,t_{k+1}))\le R_\phi$ and
$\Var_{q^{(k)}_\lambda[h_\phi]}(\log h_\phi)\in[\mu_\phi,L_\phi]$ with $0<\mu_\phi\le L_\phi\le R_\phi^2/4$;
and the diffusion-time regularity of Assumption~\ref{as:drift} holds for the $h_\phi$-dual with
constant $G_\phi$. (These are Assumptions~\ref{as:alive}--\ref{as:drift} restated for $h_\phi$.)
\end{assumption}

\paragraph{Effective residual on the common support.} Where both $q^{(k)}_{\lambda_k}[h_\phi]$
and the true Doob kernel $p^{(k)}_\star$ are supported (i.e.\ on $\Aset^\cap_k$, absent false
positives; see below), the common-base log-potential difference is
\begin{equation}
\label{eq:eta-res}
\varepsilon_k\triangleq\lambda_k\log h_\phi(\cdot,t_{k+1})-\log\hs(\cdot,t_{k+1})
=(\lambda_k-1)\log h_\phi(\cdot,t_{k+1})+\big(\log h_\phi-\log\hs\big)(\cdot,t_{k+1}).
\end{equation}

\begin{proposition}[worst-case residual at a fixed exponent]
\label{prop:eta-residual}
Under Assumptions~\ref{as:eta} and~\ref{as:eta-reg}, for $k\le N-2$, on $\Aset^\cap_k$,
\[
\osc_{\Aset^\cap_k}(\varepsilon_k)\ \le\ |\lambda_k-1|\,R_\phi\ +\ \osc_{\Aset^\cap_k}\!\big(\log h_\phi-\log\hs\big)
\ \le\ |\lambda_k-1|\,R_\phi+2\eta .
\]
This is a worst-case \emph{upper} bound, not a floor: a pure log-offset
$\log h_\phi=\log\hs+c$ has $\osc(\log h_\phi-\log\hs)=0$ and contributes nothing despite
$\eta=|c|$; what matters is the \emph{oscillation} of the log-error, not its level.
\end{proposition}
\begin{proof}
$\osc(a+b)\le\osc(a)+\osc(b)$, $\osc(c\,f)=|c|\osc(f)$, $\osc(\log h_\phi)\le R_\phi$
(Assumption~\ref{as:eta-reg}), and $\osc(\log h_\phi-\log\hs)\le2\|\log h_\phi-\log\hs\|_{\infty,\Aset^\cap_k}\le2\eta$,
applied to \eqref{eq:eta-res}.
\end{proof}

\subsection{Log-affine miscalibration: the optimal exponent is $\alpha_k\ne1$}

\begin{assumption}[no support mismatch; log-affine miscalibration]
\label{as:affine}
\textup{(a)} \emph{Common zero set:} $p^{(k)}(y\mid x)>0\Rightarrow\big(h_\phi(y,t_{k+1})>0\Leftrightarrow\hs(y,t_{k+1})>0\big)$
(so $\Aset^+_k=\varnothing$ and $\Aset^\cap_k$ is the alive support of \emph{both}).
\textup{(b)} \emph{Log-affine:} there is a step scalar $\alpha_k>0$ (uniform over predecessors)
and, for each $x_k$, a constant $c_k(x_k)\in\R$ with
$\log\hs(\cdot,t_{k+1})=\alpha_k\log h_\phi(\cdot,t_{k+1})+c_k(x_k)$ on $\Aset^\cap_k(x_k)$.
\end{assumption}

Part (a) is the (strong) statement that $h_\phi$ misranks but never mis-classifies
reachability; part (b) is the typical regression failure, a per-step mis-scaling of the
log-discriminator. Then $\hs=e^{c_k}h_\phi^{\alpha_k}$ on the alive support, and $\alpha_k>1$
means $h_\phi$ is \emph{under-confident} (flatter in log-space).

\begin{proposition}[reduction to the exact regime with rescaled exponent]
\label{prop:eta-reduction}
Under Assumptions~\ref{as:exact}(b), \ref{as:alive}, \ref{as:eta}, \ref{as:affine}, and a predictable
schedule $\lambda_k\in[\lmin,\lmax]$, the realised kernel equals the \emph{true}-$\hs$ tempered kernel at
$\tilde\lambda_k\triangleq\lambda_k/\alpha_k>0$:
\begin{equation}
\label{eq:eta-reduce}
q^{(k)}_{\lambda_k}[h_\phi](\cdot\mid x_k)=q^{(k)}_{\tilde\lambda_k}[\hs](\cdot\mid x_k).
\end{equation}
Writing $R_\star$ for the bound on $\osc(\log\hs)$ over the alive support (Assumption~\ref{as:alive}),
\textup{(i)} $P_Q(x_N\in S)=1$, and \textup{(ii)}
\[
\KL\!\big(q(x_N)\,\big\|\,p_{\text{data}}(\cdot\mid S)\big)
\ \le\ \frac{R_\star^2}{8}\sum_{k=0}^{N-2}\E_Q\Big[\big(\tfrac{\lambda_k}{\alpha_k}-1\big)^2\Big].
\]
The displayed bound is minimised by $\lambda_k=\alpha_k$ (giving $Q=P^\star$, bound $=0$) when
$\alpha_k\in[\lmin,\lmax]$ --- uniquely for steps $k\le N-2$ with $\osc_{\Aset^\cap_k}(\log\hs)>0$;
if $\alpha_k\notin[\lmin,\lmax]$ the box-clipped minimiser is the nearer endpoint. The fixed
unit exponent $\lambda_k\equiv1$ leaves $\tfrac{R_\star^2}{8}\sum_{k=0}^{N-2}\E_Q[(\alpha_k^{-1}-1)^2]$.
\end{proposition}
\begin{proof}
On $\Aset^\cap_k$, $h_\phi^{\lambda}=\exp(\lambda\log h_\phi)=\exp\!\big(\tfrac{\lambda}{\alpha_k}(\log\hs-c_k)\big)
=e^{-\lambda c_k/\alpha_k}\hs^{\lambda/\alpha_k}$; by Assumption~\ref{as:affine}(a) the two kernels
have identical support, so $p^{(k)}h_\phi^{\lambda}\propto p^{(k)}\hs^{\lambda/\alpha_k}$, which is
\eqref{eq:eta-reduce}. Thus $Q$ is the $\hs$-tempered chain with predictable exponent
$\tilde\lambda_k>0$, and Thm.~\ref{thm:sat}(b) (giving (i)) and Thm.~\ref{thm:master} (giving (ii),
with its constant $R_\star$ since the chain is $\hs$-tempered) apply verbatim. The terminal term
vanishes ($\hs\equiv1$ on the alive terminal support); the upper bound is the pointwise minimum
at $\tilde\lambda_k=1$, i.e.\ $\lambda_k=\alpha_k$, unique on steps with nonzero $\osc(\log\hs)$.
\end{proof}

So under a multiplicative log-miscalibration the \emph{right} exponent is $\alpha_k$, not $1$:
a single per-step scalar temperature \emph{exactly} corrects the discriminator (recovering
$P^\star$), and $\alpha_k>1$ forces $\lambda>1$ while $\alpha_k<1$ forces $\lambda<1$. A fixed
unit exponent cannot do this. (The crude Prop.~\ref{prop:eta-residual} hid this because it
bounded $\osc(\log h_\phi-\log\hs)$ adversarially; here that oscillation is $|1-\alpha_k|\osc(\log h_\phi)$,
removed by $\lambda_k=\alpha_k$.)

\subsection{Beyond the affine model: reachability and the satisfaction floor}

With genuine support mismatch ($\Aset^+_k\ne\varnothing$) satisfaction is no longer free, and
no single exponent restores it. We record what tempering \emph{does} control.

\begin{proposition}[reachability monotonicity and a budgeted floor]
\label{prop:eta-reach}
Under Assumption~\ref{as:eta}, for the realised kernel $q^{(k)}_\lambda[h_\phi]$ and any $\lambda>0$:
\begin{gather*}
\textup{(i)}\quad \frac{d}{d\lambda}\,\E_{q^{(k)}_\lambda[h_\phi]}[\hs]
=\Cov_{q^{(k)}_\lambda[h_\phi]}\big(\hs,\log h_\phi\big),
\end{gather*}
which is $\ge0$ iff $\hs$ and $\log h_\phi$ are positively associated under $q^{(k)}_\lambda$
(e.g.\ if $\hs$ is a nondecreasing function of $h_\phi$ on the support); under negative
association it is $<0$. \textup{(ii)} If $\E_{q^{(k)}_\lambda[h_\phi]}[-\log h_\phi]\le C$ then
\[
\E_{q^{(k)}_\lambda[h_\phi]}[\hs]\ \ge\ \big(e^{-(C+\eta)}-e^{-\eta}\,m^+_k\big)_+,
\qquad m^+_k\triangleq\E_{q^{(k)}_\lambda[h_\phi]}\big[h_\phi\,\mathbf 1\{\hs=0\}\big]
\]
the $h_\phi$-mass on false positives (the bound is informative only when the bracket is positive).
At the terminal step $k=N-1$, if the budget holds with slack $\tau$
($\E[-\log h_\phi(x_N,t_N)\mid\F_{N-1}]\le C+\tau$), then
$P_Q(x_N\in S)\ge\big(e^{-(C+\tau+\eta)}-e^{-\eta}\E_Q[m^+_{N-1}]\big)_+$.
\end{proposition}
\begin{proof}
(i) $\partial_\lambda q^{(k)}_\lambda(y)=q^{(k)}_\lambda(y)\big(\log h_\phi(y)-\E_{q^{(k)}_\lambda}[\log h_\phi]\big)$,
so $\partial_\lambda\E_{q^{(k)}_\lambda}[\hs]=\Cov_{q^{(k)}_\lambda}(\hs,\log h_\phi)$.
(ii) Jensen on the convex $\exp$: $\E_{q^{(k)}_\lambda}[h_\phi]\ge e^{\E[\log h_\phi]}=e^{-\E[-\log h_\phi]}\ge e^{-C}$.
Split the support into $\Aset^\cap_k$ ($\hs>0$) and $\Aset^+_k$ ($\hs=0$); on $\Aset^\cap_k$,
$\hs\ge h_\phi e^{-\eta}$ (Assumption~\ref{as:eta}), so
$\E[\hs]=\E[\hs\mathbf 1_{\Aset^\cap_k}]\ge e^{-\eta}\E[h_\phi\mathbf 1_{\Aset^\cap_k}]
=e^{-\eta}(\E[h_\phi]-m^+_k)\ge e^{-(C+\eta)}-e^{-\eta}m^+_k$; since $\E[\hs]\ge0$ the positive part follows.
At $k=N-1$, $\hs(\cdot,t_N)=\mathbf 1\{\cdot\in S\}$ gives $\E_Q[\hs(x_N,t_N)]=P_Q(x_N\in S)$;
the conditional budget with slack $\tau$ replaces $C$ by $C+\tau$.
\end{proof}

A tighter budget $C$ (a larger $\lambda$, when feasible and not box-clipped) raises the
$e^{-(C+\eta)}$ term; but the false-positive mass $m^+_k$ subtracts an amount not removable by
tempering alone --- the satisfaction loss of an over-permissive $h_\phi$, which needs a better
$h_\phi$ or the hard $\epsilon$-floor.

\subsection{The dual ascent finds the corrective exponent}

\begin{remark}[the tracking machinery transfers, with $h_\phi$-constants]
\label{rem:eta-bridge}
Lemmas~\ref{lem:noise}--\ref{lem:contr} and Theorem~\ref{thm:track} used only that the tilt
function is bounded in $[\epsilon,1]$, with a nondegenerate concave dual and time-regular
gradient; they hold verbatim for $h_\phi$ once Assumption~\ref{as:eta-reg} supplies the
$h_\phi$-constants $(R_\phi,\mu_\phi,L_\phi,G_\phi)$ in place of $(R,\mu,L,G)$. Then $\lambda_k$
tracks the $h_\phi$-dual optimum $\lambda^\star_k$ solving $\E_{q^{(k)}_\lambda[h_\phi]}[-\log h_\phi]=C$.
Under Assumption~\ref{as:affine}, if $\alpha_k\in[\lmin,\lmax]$ then the matched budget
$C^{\phi\star}_k\triangleq\E_{q^{(k)}_{\alpha_k}[h_\phi]}[-\log h_\phi]$ gives $\lambda^\star_k=\alpha_k$,
and the $h_\phi$-analogue of Lemma~\ref{lem:calib} (using $S^\phi_k{}'=-\Var_{q[h_\phi]}(\log h_\phi)\in[-L_\phi,-\mu_\phi]$)
gives the calibration gap $|\lambda^\star_k-\alpha_k|\le|C-C^{\phi\star}_k|/\mu_\phi$. \emph{Thus
for a general fixed $C$ the adaptive dual tracks the $h_\phi$-budget optimum $\lambda^\star_k(C)$,
and under a matched (or near-matched) budget it is an online estimator of the
miscalibration-correcting exponent $\alpha_k$}, enforcing a reachability budget from the only
object it has, $h_\phi$.
\end{remark}

\begin{remark}[positioning and scope]
\label{rem:eta-position}
\emph{Theory.} The optimal exponent under log-affine miscalibration is $\alpha_k$, generally
$\ne1$ and step-dependent (Prop.~\ref{prop:eta-reduction}); a fixed unit exponent is therefore
suboptimal, and adadual tracks the $h_\phi$-budget optimum $\lambda^\star_k(C)$ --- equal to
$\alpha_k$ under a matched budget (Rem.~\ref{rem:eta-bridge}): $\alpha_k>1$ (under-confident
$h_\phi$) calls for stronger tempering, $\alpha_k<1$ for weaker. The
exact-discriminator sections are the special case $\alpha_k\equiv1$, $\eta=0$, where $\lambda\equiv1$
is optimal and adadual is unnecessary --- the value of adaptivity is exactly an $\eta>0$
phenomenon. \emph{Two effects of $\lambda$.} Within the common support, $\lambda_k$ corrects
calibration (exactly at $\alpha_k$); across the support boundary, raising $\lambda_k$ raises true
reachability only under positive association (Prop.~\ref{prop:eta-reach}(i)) and lifts the
budgeted floor (ii). \emph{Limits.} The false-positive mass $m^+_k$ (support mismatch) caps
satisfaction and is not guaranteed removable by any single scalar exponent --- it needs a better
$h_\phi$ or the hard $\epsilon$-floor. \emph{Empirically} (motivation, outside the theorem chain),
the learned-unit baseline $\lambda\equiv1$ plateaus near $45\%$ on QM9 novelty, consistent with
$\alpha_k>1$ and nonzero $m^+_k$.
\end{remark}

\subsection{Over-solving the inner dual: an interior-optimal inner-iteration count $J$}
\label{sec:eta-oversolve}

The frozen-step solver ascends the concave $h_\phi$-dual from $\lambda=0$; after $J$ inner
iterations the exponent $\lambda_J$ increases monotonically to the budget optimum
$\lambda^\star_k(C)$ solving $\E_{q^{(k)}_\lambda[h_\phi]}[-\log h_\phi]=C$
(Rem.~\ref{rem:eta-bridge}). So a larger inner budget $J$ means a \emph{larger} realised
exponent. We show that the per-step \emph{true} satisfaction
\begin{equation}
\label{eq:Vdef}
V(\lambda)\ \triangleq\ \E_{q^{(k)}_\lambda[h_\phi]}\!\big[\hs(\cdot,t_{k+1})\big]
\qquad(\text{per-step violation }=1-V)
\end{equation}
is, under an imperfect discriminator, \emph{non-monotone} in $\lambda$ --- hence in $J$ ---
so solving the dual to convergence overshoots the satisfaction-optimal exponent. Throughout
write $T\triangleq\log h_\phi(\cdot,t_{k+1})$, so $q^{(k)}_\lambda[h_\phi]$ is the exponential
family $\propto p^{(k)}(\cdot\mid x_k)\,e^{\lambda T}$ with natural parameter $\lambda$ and
sufficient statistic $T$, and (Prop.~\ref{prop:eta-reach}(i)) $V'(\lambda)=\Cov_{q^{(k)}_\lambda[h_\phi]}(\hs,T)$.

\begin{corollary}[well-specified ranking $\Rightarrow$ monotone $\Rightarrow$ over-solving is safe]
\label{cor:comonotone}
Fix $k\le N-2$ and $x_k$. If $\hs(\cdot,t_{k+1})$ is a nondecreasing function of $T$ on the
realised support (\emph{comonotone}; in particular the exact case $h_\phi=\hs$, where
$\hs=e^{T}$), then
\[
V'(\lambda)=\Cov_{q^{(k)}_\lambda[h_\phi]}(\hs,T)\ \ge\ 0\qquad\forall\,\lambda\ge0,
\]
so $V$ is nondecreasing and a larger inner budget $J$ never lowers per-step satisfaction.
Consequently a \emph{necessary} condition for ``larger $J$ strictly lowers $V$'' is a rank
inversion of $h_\phi$ relative to $\hs$: a pair $y,y'$ in the support with $T(y)>T(y')$ yet
$\hs(y,t_{k+1})<\hs(y',t_{k+1})$.
\end{corollary}
\begin{proof}
$V'(\lambda)=\Cov_{q^{(k)}_\lambda[h_\phi]}(\hs,T)$ by Prop.~\ref{prop:eta-reach}(i). For a single
random variable $T$ and a nondecreasing $g$, $\Cov(g(T),T)\ge0$ (Chebyshev's association
inequality); apply with $\hs=g(T)$. The exact case has $T=\log\hs$ and $g=\exp$. Necessity is
the contrapositive: absent any rank inversion, $\hs$ is a nondecreasing function of $T$.
\end{proof}

\begin{proposition}[interior-optimal exponent and strict over-solving penalty]
\label{prop:eta-oversolve}
Fix $k\le N-2$ and $x_k$; let $\mathcal Y\triangleq\{y:p^{(k)}(y\mid x_k)>0\}$ and
$\mathcal Y^\star\triangleq\argmax_{y\in\mathcal Y}T(y)$, with $T$ nonconstant on $\mathcal Y$.
Assume
\begin{itemize}
\item[\textup{(A1)}] \emph{initial improvement:} $\Cov_{q^{(k)}_0[h_\phi]}(\hs,T)>0$ (the
discriminator ranks correctly under the base kernel);
\item[\textup{(A2)}] \emph{mode suboptimality / false positives:}
$\displaystyle\lim_{\lambda\to\infty}V(\lambda)=\E_{q^{(k)}_\infty}[\hs]\le V(0)$, where
$q^{(k)}_\infty$ is $p^{(k)}(\cdot\mid x_k)$ restricted to $\mathcal Y^\star$.
\end{itemize}
Then $\argmax_{\lambda\ge0}V\subseteq(0,\infty)$. Let $\lambda^\dagger$ be its largest element
and $C^\dagger\triangleq\E_{q^{(k)}_{\lambda^\dagger}[h_\phi]}[-T]$. For every budget
$C<C^\dagger$ the converged exponent satisfies $\lambda^\star_k(C)>\lambda^\dagger$ and
\[
V\big(\lambda^\star_k(C)\big)\ <\ V(\lambda^\dagger)\ =\ \max_{\lambda\ge0}V(\lambda).
\]
\end{proposition}
\begin{proof}
$V$ is real-analytic in $\lambda$ (exponential family). (A1) gives $V'(0)>0$, hence
$\sup_\lambda V>V(0)$. As $\lambda\to\infty$ the tilt concentrates on $\mathcal Y^\star$ (Gibbs
concentration on the finite set $\mathcal Y$), so $V(\lambda)\to\E_{q^{(k)}_\infty}[\hs]\le
V(0)<\sup V$ by (A2). A continuous function on the compactified $[0,\infty]$ whose values at
both endpoints lie strictly below its supremum attains that supremum at an interior point, so
$\lambda^\dagger\in(0,\infty)$. The budget map $B(\lambda)\triangleq\E_{q^{(k)}_\lambda[h_\phi]}[-T]$
has $B'(\lambda)=-\Var_{q^{(k)}_\lambda[h_\phi]}(T)<0$ (strict, as $T$ is nonconstant), so $B$ is
a decreasing bijection and $C<C^\dagger\Rightarrow\lambda^\star_k(C)=B^{-1}(C)>B^{-1}(C^\dagger)=\lambda^\dagger$.
As $\lambda^\dagger$ is the largest maximiser, $V(\lambda)<\max V$ for all $\lambda>\lambda^\dagger$,
in particular at $\lambda^\star_k(C)$.
\end{proof}

\begin{corollary}[interior-optimal inner budget; early stopping is implicit budget relaxation]
\label{cor:Jdagger}
Run the inner ascent from $\lambda_0=0$ with steps small enough that the iterates
$\{\lambda_J\}_{J\ge0}$ increase monotonically to $\lambda^\star_k(C)$ (projected gradient ascent
on the concave dual). Under (A1)--(A2) with $C<C^\dagger$ there is a finite $J^\dagger$ with
$\lambda_{J^\dagger}\approx\lambda^\dagger$ minimising per-step violation $1-V$; for $J>J^\dagger$
the exponent overshoots $\lambda^\dagger$ and per-step violation strictly increases. Equivalently,
stopping at $J^\dagger$ realises the looser, satisfaction-optimal budget $C^\dagger>C$ rather than
the configured $C$: \emph{capping the inner iterations is implicit budget relaxation.}
\end{corollary}

\begin{remark}[reading and scope]
\label{rem:oversolve-scope}
Together Cor.~\ref{cor:comonotone} and Prop.~\ref{prop:eta-oversolve} explain the empirical
``over-solving hurts'': driving the inner loop toward the converged $\lambda^\star_k(C)$ exploits
$h_\phi$'s rank inversions and false positives (Goodhart on the proxy), so true satisfaction $V$
is single-humped in $\lambda$ with an interior peak $\lambda^\dagger$ and an interior inner budget
$J^\dagger$ is optimal. Under a correctly-ranking discriminator ($\eta=0$, Cor.~\ref{cor:comonotone})
$V$ is monotone and no finite $J$ is preferred --- so ``larger $J$ raises violation'' is itself an
$\eta>0$ signature, consistent with the QM9 plateau (Rem.~\ref{rem:eta-position}). Two caveats.
\textup{(i)} $V$ is the \emph{per-step} surrogate $\E_{q_\lambda}[\hs]$; terminal violation
telescopes the per-step martingale, so the single-hump statement is qualitative at the trajectory
level. \textup{(ii)} (A2) uses the sufficient form $V(\infty)\le V(0)$; the conclusion holds under
the weaker $V(\infty)<\sup_\lambda V$.
\end{remark}

\end{document}
